\uuid{fvUO}
\titre{Congruences}
\auteur{bodin ; croizat ; sacre}
\organisation{exo7}
\licence{CC BY-NC-SA 4.0}
\datecreate{2022-03-08}
\niveau{L1}
\module{Algèbre}
\chapitre{Arithmétique dans Z}
\sousChapitre{Autre}
\difficulte{2}
\theme{Congruences}
\exo7qcmid{hilisit.arithmetique.047}

\contenu{
\begin{question}
\qclassification{103.99}
    Soit $k$ un entier et $N = 5k^2-10k+4$. On peut affirmer :
    \begin{answers}
        \good{$N \equiv 4 \;[5]$}
        \bad{$N \equiv 5 \; [5]$}
        \good{$N \equiv 5k^2 \;[2]$}
        \bad{$N \equiv 1 \;[2]$}
    \end{answers}
    \begin{explanations}
        Puisque $5 \equiv 10 \equiv 0 \;[5]$, on a $5k^2-10k \equiv 0 \;[5]$. Donc $N \equiv 4 \;[5]$.\\
        D'autre part, puisque $10 \equiv 4 \equiv 0 \;[2]$, on a $-10k+4 \equiv 0 \;[2]$. Donc $N \equiv 5k^2 \;[2]$. Le cas $k=2$ (ou tout autre entier pair) montre que l'on peut avoir $N \equiv 5k^2 \equiv 0 \;[2]$.
    \end{explanations}
\end{question}
}
